% Source: https://artofproblemsolving.com/community/c3h1402576_algebra_manipulation
% Stable identity: aops-topic-1402576-q01
\textbf{$d=0$.} For each $x>0$ the right side $2x^3+3x^2+x=x(x+1)(2x+1)$ is positive, so both $(x,y)$ and $(x,-y)$ lie on the curve for some $y\ne0$. They have the same $x_1$, so $(cy+d)^2=(-cy+d)^2$, that is $4cdy=0$, and since $c\ne0$, $d=0$.

\textbf{The other constants.} Now $y_1^2=c^2y^2=\dfrac{c^2}{6}(2x^3+3x^2+x)$ must equal $(ax+b)^3-36(ax+b)$ for all $x>0$, so the two cubics in $x$ are identical. Comparing coefficients:
\[
x^3:\ a^3=\dfrac{c^2}{3},\qquad x^2:\ 3a^2b=\dfrac{c^2}{2},\qquad x:\ 3ab^2-36a=\dfrac{c^2}{6},\qquad 1:\ b^3-36b=0 .
\]
The first two give $3a^2b=\dfrac{3a^3}{2}$, so $b=\dfrac{a}{2}$. Then the third gives $\dfrac{3a^3}{4}-36a=\dfrac{a^3}{2}$, so $a^2=144$; and $a>0$ because $a^3=\dfrac{c^2}{3}>0$; so $a=12$, $b=6$, and $c^2=3a^3=5184$, $c=72$. The constant term checks: $6^3-36\cdot6=0$.

Indeed $x_1=12x+6=6(2x+1)$ gives $x_1^3-36x_1=216(2x+1)\left((2x+1)^2-1\right)=864\,x(x+1)(2x+1)=5184y^2=(72y)^2$. Hence
\[
a+b+c+d=12+6+72+0=\boxed{90}.
\]
